%!TEX TS-program = pdflatex
%!TEX encoding = UTF-8 Unicode
%!TEX spellcheck = English (Aspell)

\documentclass[11pt]{article}
\usepackage{../../jeffe,../../handout,graphicx}
\usepackage{braket}


% =========================================================
\begin{document}

Last class, we learned that quantum operations must satisfy
the following properties:
\begin{itemize}
    \item \textbf{Unitary}: Given a superposition of $n$ qubits 
    \[\ket{\psi} = \sum_{x \in \{0,1\}^n} \alpha_x \ket{x},\]
    where $\alpha_x$ is the amplitude of the basic state $\ket{x}$
    so that $\sum_{x \in \{0,1\}^n} \alpha_x^2 = 1$, a quantum operation $U$ must be unitary,
    meaning that after applying $U$ to $\ket{\psi}$, it is still true that $\sum_{x \in \{0,1\}^n} \alpha_x^2 = 1$.
    \item \textbf{Reversible}: A quantum operation $U$ must be reversible,
    meaning that there must exist some quantum operation $U^{-1}$ such that
    applying $U$ and then $U^{-1}$ to any superposition of $n$ qubits
    $\ket{\psi}$ results in the original state $\ket{\psi}$.
\end{itemize}

For the following operations, either give their reversal operations or indicate that
they are not reversible. We use capital letters to denote individual qubits; for example,
if we have two qubits, we refer to the first qubit as $A$ and the second qubit as $B$.

\begin{table}[h]
\begin{tabular}{c|c|c}
    Operation $U$ & Reversal Operation $U^{-1}$ & Minimal example\\
    \hline
    $\neg A$ & $\neg A$ & $\ket{0} \to \ket{1} \to \ket{0}$ \\
     &  & $\ket{1} \to \ket{0} \to \ket{1}$ \\
     \hline
    if $A$ then $\neg$ B (returns both qubits) &   \\
    (e.g., $\ket{00} \to \ket{00}$) & & \\
    & & \\
    & & \\
    & & \\
    & & \\
    & & \\
    \hline
    $A \oplus B$ &   \\
    returns only one qubit& & \\
    (e.g., $\ket{00} \to \ket{0}$) & & \\
    & & \\
    & & \\
    & & \\
    & & \\
    \hline
    Hadamard $A$, defined as: & & \\
    $\ket{0} \to \sqrt{\frac{1}{2}}$ amplitude on $\ket{0}$ & & \\
    and $\sqrt{\frac{1}{2}}$ amplitude on $\ket{1}$ & & \\
    $\ket{1} \to \sqrt{\frac{1}{2}}$ amplitude on $\ket{0}$ & & \\
    and $-\sqrt{\frac{1}{2}}$ amplitude on $\ket{1}$ & & \\
\end{tabular}
\end{table}

\end{document}